\section{One-mark plateau estimate by Hilbert--Schmidt links}\label{c4:plateau}

The determinant calculation has established that the homogeneous sum
\(\tau_N\) does not vanish. We now estimate the marked sum relative to it.

\begin{proposition}[The homogeneous one-mark ratio]\label{prop:one-mark-ratio}
For the continued homogeneous color sums of
Proposition~\ref{a:finite-polygons},
\[
                    \big|\tau^v_{N-1}/\tau_N\big|=N^{-2/3+o(1)}.
\]
\end{proposition}

The proof uses an independent positive lower bound. Equations
\eqref{c4:Press-3}, \eqref{c4:Press-6} and \eqref{a:finite-winding} give,
for some \(c>0\) and all sufficiently large \(N\),
\begin{equation}\label{eq:one-mark-positive-input}
 \big|\tau^v_{N-1}/\tau_N\big|=E_N/c_*\ge A_N/c_*\ge c/N.
\end{equation}
We will use this bound after the radial calculation to show that its
leading coefficient cannot vanish.

Continue the radial notation with \(\epsilon=1,\ a=7/18\), and work first
with each summand of the rotated compensator. By the results on \(F_b\),
removing \(Nh\) from the central charge \(B\) extracts the self factor
\(d_R(h)^{N^2}\) and a unit phase common to all summands and configurations.
Its interactions with its own images and with the compensator tend to 1
in the successive limits at fixed \(N\). Its remaining effect is to
multiply the root activities by
\[
 \mathcal B_b(u)=\exp(N F_b(u+i\phi_b))
\]
including periodic images of \(F_b\) for finite \(L\). The remaining charge at 0 is \(p\). For strict types \(b\ne\beta\), \(|\mathcal B_b|\le\exp(-cN e^{-|u|/6})\), also with images; for \(\beta\) it is 1. Factoring the compensator interactions as above is harmless provided the remaining amplitudes per summand stay bounded through the successive limits at fixed \(N\), proved below. Also \(Z_C^r(L)\to1\). We construct operators for the remaining \(\exp(7L/192)Z_R^r(L)\) times the radial gas.

\subsection{Hilbert--Schmidt links}

Use boxes \([jD,(j+1)D)\), \(D\) large, \(L=M D,\ u_*=kD\), \(k\) and \(M\) integers with both marks well inside the period. Configurations \(x_j\) in a box are their root vertices (angles/types with relative positions \(\xi\in[0,D)\)), symmetric measure per type \(\prod d\xi/(2\pi)\) divided by the count factorial. Also append deterministic marks where present (at \(\xi=0\)) when computing interactions/weights. Entrance slope \(\ell_j\) is before these, in \(\mathcal M_t=a+t+\mathbb Z/p^2\) where \(t=0\) or 1 according to being outside or inside the marked interval at entrance (\(j=0\) still outside, \(j=k\) still inside). Impose across a link \(j\to j+1\)
\[
 \ell_{j+1}=\ell_j+\sum_{b\in x_j}^{\rm incl.\ marks} n_b/p^2 .
\]
Use states \(L^2(x_j,\ell_j;\mathscr F)\), where \(\mathscr F\) is a positive-metric bosonic Fock Hilbert space: direct sum of symmetric powers of the direct sum of the positive-frequency blocks of \eqref{c4:Rad-2}, using fixed ordinary positive inner products (same for repetitions of each block). Write \(\Gamma_{\rm op}(P')\) for second quantization, and \(a^\dagger(F),a(G)\) for standard normalized creation/annihilation with vector and linear covector, \([a(G),a^\dagger(F)]=G(F)\). Identify repeated sector state spaces by relative box coordinates.

Each link is a kernel acting from the old/source box \(j\) to the
new/target \(j+1\), zero without compatibility. We assign within-box
interactions to the two links meeting that box, and adjacent-box
interactions to their common link. These scalar factors provide the
local count bound. The damped Fock factor below will supply every
interaction between boxes at distance at least two. Activities and
local phases belong to the source box. More precisely, the scalar
part includes:
\par\noindent\textbullet\quad \(\exp(7D/192)\), heat cost over the source box, local factors \eqref{c4:Rad-1} and all residual phases therein, contour coefficients and any profile multipliers in that source (not the target).
\par\noindent\textbullet\quad The exponential of minus short interactions \(\mathcal Q\) as follows. Choose smooth weights \(f(\xi/D),g(\xi/D)\) in each box, both \(\ge\eta>0\) small, \(f^2+g^2=1\), \(f=\eta\) on the leftmost quarter and \(g=\eta\) on the rightmost quarter. Within the old box include each \(x<y\) with fraction \(f_x f_y+(1-f_x f_y-g_x g_y)/2\); within the new the analogous fraction with \(g\) instead of the first \(f\)-product. Across old to new include all forward pairs with full weight.

Additionally the kernel contains the operator
\[
 e^{a^\dagger(\mathfrak F_j)}\Gamma_{\rm op}(\mathfrak P_D)e^{a(\mathfrak G_{j+1})},\qquad
 \mathfrak P_D=e^{-D\Lambda}\ {\rm blockwise},
\]
where including marks
\[
 \mathfrak F_j|_l=\sum_{x\in j}e^{-\Lambda_l(3D/2-\xi_x-i\phi_x)}P_l b_x,\qquad
 \mathfrak G_j|_l=-\sum_{x\in j}\big(e^{-\Lambda_l(\xi_x+D/2+i\phi_x)}\Lambda_l^{-1}(\cdot),b_x\big).
\]

These links are Hilbert--Schmidt, bounded locally uniformly for
sufficiently large widths \(D\) and bounded profiles, and Hilbert--Schmidt
Lipschitz in uniform changes of bounded source profiles. To see this,
apply \eqref{c4:Rad-3} first with weights \(f\) on the old box and \(g\)
on the new box. The within-box discrepancy from this weighting to the
actual fractions is
\[
 \frac{1-f_xf_y-g_xg_y}{2}
 =\frac14\big\|(f_x,g_x)-(f_y,g_y)\big\|^2
 =O\big((|u_x-u_y|/D)^2\big).
\]
It is nonnegative. Negative parts of the kernels are bounded with
exponential decay, while identical-root singularities are favorable,
so this discrepancy costs at most \(C D^{-2}\) times ordinary
squared-cell sums. Across the link the discrepancy \(1-f_xg_y\) is
\(\eta^2\) if both vertices are in adjoining outer quarters. There
\(f_x\) and \(g_y\) are close to 1, so their weighted counts control
the unweighted counts uniformly; elsewhere the real separation is at
least \(D/4\). Choosing \(\eta\) small and then \(D\) large therefore
retains a short-interaction gap
\(c\sum_I n_I^2-C(n_{\rm tot}+1)\), with \(c>0\) independent of large
widths. Here \(n_I\) are unit-cell counts and \(n_{\rm tot}\) is the total
count in the two boxes. Marks on distinct rays cost at most linearly
as well. The Fock part has Hilbert--Schmidt norm
\[
 \le C\exp(C(\|\mathfrak F_j\|^2+\|\mathfrak G_{j+1}\|^2)),\qquad
 \|\mathfrak F_j\|+\|\mathfrak G_{j+1}\|\le C e^{-c' D}(n_{\rm tot}+1).
\]
In fact split \(\Gamma_{\rm op}(\mathfrak P_D)\) into three equal dampings; the middle is uniformly Hilbert--Schmidt (sum squared singular values via the product over one-particle singular values, using small operator and Hilbert--Schmidt norms there by mode decay). Each outer exponential times damping is operator bounded as stated: \((a^\dagger(F))^{j'}\) on degree \(n\) has norm \(\le\|F\|^{j'}\sqrt{(n+j')!/n!}\), giving the bound after degree damping of at least \(2^{-n}\), and adjointly for annihilators. Thus the Fock cost is absorbed by large \(D\) (only \(O(D+1)\) unit cells). The source heat, summed over entrance slopes, costs only \(O(n_{\rm tot}+1)\), since slope stays within \(O(n_{\rm tot}+1)\) of entrance, and exit is tied by compatibility. These estimates integrate even with fixed exponential factors in counts, yielding the Hilbert--Schmidt and difference assertions by absolute bounds and product telescoping. Kernels with the fixed/limiting profiles used below also depend continuously on \(D\) in this norm (isometrically use \([0,1)\) relative coordinates; pointwise continuity off null sets including for the Fock part by the damping estimates, and the same domination locally).

Finite-mode Bose trace--determinant formulas are discussed in \cite[Section~2]{Klich2002}. Here the linear source insertions and the passage to infinitely many modes also require the damping estimates just established.

Tracing the cyclic link product (say \(M\ge3\)) now gives the desired \(\exp(7L/192)Z_R^r\) times the gas. Scalar interactions account for same and adjacent boxes on the extension. The Fock trace produces contractions \(\mathfrak G_{j'}\mathfrak P_D^{j'-j-2}\mathfrak F_j\) for all \(j'\ge j+2\) on that extension, precisely the remaining minus \(\mathcal Q\) terms, and the determinant \(\det(1-\mathfrak P_D^M)^{-1}=Z_R^r\). For verification, commuting into \(e^{a^\dagger(F)}\Gamma_{\rm op}(\mathfrak P_D^M)e^{a(G)}\) gives the internal contractions, and its trace adds \(G(1-\mathfrak P_D^M)^{-1}F\) in the exponential, where \(F=\sum_{r=0}^{M-1}\mathfrak P_D^{M-1-r}\mathfrak F_{j+r}\), \(G=\sum_{r=1}^M \mathfrak G_{j+r}\mathfrak P_D^{r-1}\). This follows at finitely many modes by the same trace identity as earlier (equivalently commute Taylor coefficients at zero and use entire analyticity by damping). Passage to all modes follows by tensor factorization and the damping bounds (on fixed box variables the tail factors differ from the tail vacuum projector by vanishing Hilbert--Schmidt norm). Kernel products and traces here are cyclic integration, valid by finite-kernel approximation since products of the fiber Hilbert--Schmidt norms integrate absolutely around the cycle bounded by the product of full Hilbert--Schmidt norms (Cauchy--Schwarz for squared-integrable kernels). The factorial measures per box reconstitute those of the full period by indistinguishability.

\subsection{Spectra from exact traces}

The link construction has recovered the full radial gas. We next extract the spectra of its stationary and suppressed-root regions from their exact trace powers. This step uses compactness, not positivity of the complex kernels.

Define baseline links \(K_t\) on each single sector by omitting all marks and setting every root profile to 1, using root character \(\exp(-2\pi i t\zeta_b)\) (these are baseline integrands for \(\mu\) shifted by \(tp\)). Define \(J_t\) similarly but with strict-root source activities zero. Both operators use the same state space: target configurations are retained, but slope compatibility depends only on the source roots, so \(J_t\) cannot decrease the entrance slope. These are complex operators. The two angular identities give
\[
 \operatorname{Tr}(K_t^M)=Z_C^r(MD)\qquad (M\ge3).
\]
Thus both \(K_t\) have simple leading eigenvalue 1, remaining spectrum in \(|\lambda|\le e^{-D/8}\), and powers tend exponentially to a rank-one projection \(\Pi_t^{\rm op}\). Also, with multiplicities,
\begin{equation}
 \operatorname{Spec}_{\ne0}(J_t)
 =\{\exp(D(7/192-p^2\ell^2/2-e)):\ell\in\mathcal M_t,\
                    e\ {\rm oscillator\ excitation\ of}\ R\}.                    \label{c4:Link-1}
\end{equation}
Here \(e\) sums positive frequencies (bosonic multiplicities, also vacuum 0); least positive \(e=1/6\). Indeed trace cycles for \(J_t\) have only \(\beta\) allowed, necessarily absent by slope increase, so powers have free \(Z_R^r\) and fixed-slope heat. The same formula on invariant flag subspaces supported on \(\ell\ge s'\) (\(J_t\) never decreases slope) and their quotients restricts the slope labels accordingly by the same traces. Consequently the Riesz range at nonzero \(\lambda\) in \eqref{c4:Link-1} is supported on slopes at least the minimum \(\ell\) for that \(\lambda\), while the Riesz projection annihilates states supported strictly above the maximum such \(\ell\) (pass to quotient or restriction).

Here extraction of these spectral conclusions from trace powers uses just compactness and resolvent integrals. Outside a circle of small regular radius there are finitely many spectral points with finite-dimensional Riesz parts, and powers after removing these decay in norm at that radius. Indeed approximate the Hilbert--Schmidt operator by finite rank plus remainder of norm \(<\) half the radius; using the remainder's resolvent reduces inversion to a finite matrix determinant, giving a meromorphic resolvent outside with finite-rank principal parts. Cauchy projections around poles split off the finite-dimensional parts (projection and commutation by the resolvent identity), and deforming the large-circle power integral gives the decay. For traces the remaining error is controlled likewise using two extra Hilbert--Schmidt factors. Thus comparison with the given absolutely convergent power lists identifies the finite parts (with algebraic multiplicities) outside each regular radius, since a mismatch of distinct powers above the radius cannot cancel to that error (Vandermonde on consecutive exponents). This argument applies also to the restrictions and quotient compressions above, so isolating contours may avoid all these spectra and respect the flag. Power estimates and finite spectral cuts are uniform/continuous locally in \(D\) by the same contours. In particular \(K_t\) and \(J_t\) are power bounded, for \(J_t\) even decaying since \(\min_{\ell\in\mathcal M_t}p^2\ell^2/2>7/192\).

\subsection{Removing the distant compensator}

Take the successive limits at each fixed \(N,D\) along the indicated box choices (the confluent-trace limits already exist). Define \(C_*(D)\), mapping sector 1 to 0, as the product link \(k\) after link \(k-1\), with the compensator at the start of box \(k\) and all root profiles equal to 1 in these two boxes, summing its finitely many possibilities with stated coefficients/phases. It is independent of \(k,N\). Define \(C_N(D)\) similarly from links \(0,-1\) with residual mark \(p\) and profiles \(\mathcal B_b\) (without images). The slope and phase rules in these links are literally the radial ones, the boundary switching to marked at \(p\) and back to exterior at the compensator. Away from marks let \(L_j(\mathcal B)\) denote the single-sector link with source profiles \(\mathcal B_b\) evaluated in box \(j\), sector 0 for \(j\le -2\) and sector 1 for \(j\ge1\) after removing the compensator to \(+\infty\).
Then the marked limiting amplitude divided by \(d_R(h)^{N^2}\), up to the common unit phases, is
\[
 \operatorname{Tr}[C_* Y^+ C_N Y^-],\quad
 Y^-=\lim_{r\to\infty}L_{-2}(\mathcal B)\cdots L_{-r}(\mathcal B)\Pi_0^{\rm op},\qquad
 Y^+=\lim_{r\to\infty}\Pi_1^{\rm op} L_r(\mathcal B)\cdots L_1(\mathcal B).
\]
Indeed send period endpoints to \(\pm\infty\) first at fixed \(k\). Profile images can then be dropped with vanishing error, their multiplier difference summed in sup norms per box bounded by \(C_D N\) times a decreasing exponential in the nearest endpoint distance to 0 (\(\Re F_b\le0\), no amplification). Profiles without images are summably close to 1 outside at fixed \(N\). Products of \(K_t+\mathcal E_j\) on any interval are bounded by \(C\exp(C\sum\|\mathcal E_j\|)\) by power boundedness and subset expansion. Thus outer tails can be replaced by powers with uniformly small error, using the same estimates on each sector segment and bounded connecting links (for cyclic traces keep two Hilbert--Schmidt links for trace norm control, splitting intervals a bounded number of times per telescoping error term). With replacements the long stationary exterior stretch tends to its projection, and undoing truncations proves convergence; bounds are also uniform for large \(k\).

Now send \(k\to\infty\). Compensation links can use the profiles 1 with vanishing error, extracted \(Nh\)-to-compensator interactions tend to 1 for each summand, and long stretches insert the exterior and marked projections as displayed, again first truncating tails. One-sided products starting at projections converge by summable perturbations. Translation of a link away from the distinguished vertices has no other effects. The normalization used leaves stationary eigenvalues exactly 1.

\subsection{The long central interval and nonvanishing}

We have reduced the marked amplitude to two one-sided products joined by fixed mark operators. The remaining calculation isolates their long regions of suppressed root activity and compares the permitted slope pairs.

For large \(N\) choose \(D\) locally near a given admissible \(D_*\) such that \(mD=6\log N\), \(m\) a positive integer. Put for strict types
\[
 G_b^-(v')=\exp(c_b^- e^{v'/6}),\qquad G_b^+(v')=\exp(c_b^+ e^{-v'/6}),
\]
and for \(\beta\) take both identically 1. On \(j\le-2\) replace \(\mathcal B_b\) by \(G_b^-(u+mD)\); on \(j\ge1\) by \(G_b^+(u-mD)\). The errors summed in sup norms per box are \(O(N^{-4})\) on either half-line, locally uniform around \(D_*\); e.g. \(v'=u+mD\) on the left gives
\[
 |\mathcal B_b(u)-G_b^-(v')|\le C N^{-4} e^{5v'/6}\exp(-c e^{v'/6})
\]
by the mode tail and negative real bounds for both exponents. In boxes \(-1,0\) set all strict-root profiles to zero with exponentially small error, writing \(C^\circ(D)\) for that central product. Total trace error after all replacements is \(O(N^{-4})\): split the products into outer pieces (relative to \(\pm mD\)) summably close to \(K_t\) powers, inner pieces summably close to \(J_t\) powers, and central links, with perturbation sums uniformly bounded by the exponential outer and superexponential inner profile tails. Then use the preceding telescope/power bounds (rank or Hilbert--Schmidt products remain in the traces). \(C^\circ\) only increases slope by a value \(\ge1\), since its source boxes now allow at most \(\beta\) roots and the one \(p\) mark.

After replacement define \(L_j^-\) as the sector-0 link with source profile \(G_b^-(jD+\xi)\), no marks; \(L_j^+\) likewise on sector 1 with \(G_b^+(-jD+\xi)\). For each side these converge exponentially at \(j\to-\infty\) to \(K_t\) and superexponentially at \(j\to+\infty\) to \(J_t\), uniformly locally in \(D\). Put
\[
 Y_n^-=\lim_{r\to\infty} L_{n-1}^-\cdots L_{-r}^-\Pi_0^{\rm op},
 \qquad Y_n^+=\lim_{r\to\infty}\Pi_1^{\rm op} L_{-r}^+\cdots L_{n-1}^+ .
\]
The desired trace now is \(\operatorname{Tr}[C_* Y_m^+ C^\circ Y_{m-1}^-]\).
For a finite spectral cut retaining \(\lambda\) in the respective \eqref{c4:Link-1} above a sufficiently small regular radius \(r_0\), locally uniform near \(D_*\),
\[
 Y_n^-=\sum_\lambda J_0^n H_\lambda^-+O(r_0^n),\qquad
 Y_n^+=\sum_\lambda H_\lambda^+ J_1^n+O(r_0^n)
\]
in norm, where \(H_\lambda^-\) has range in the Riesz range and \(H_\lambda^+\) includes the corresponding projection on the right; these coefficients are continuous in \(D\). Indeed \(Y_{n+1}^-=J_0Y_n^-+\mathcal E_n Y_n^-\) with superexponentially small \(\mathcal E_n\) and uniformly bounded \(Y_n^-\). On retained Riesz parts propagate corrections backward by inverse powers (uniformly convergent, tail superexponentially small); on the complement convolve forward using the radius-cut power estimate. The right-factor version is identical.

Let \(P_{0,L}\) and \(P_{1,R}\) be a retained pair of Riesz projections
in sectors 0 and 1. The corresponding term contains
\(P_{1,R}C^\circ P_{0,L}\). The range of \(P_{0,L}\) is supported on
slopes at least its smallest spectral label. The central operator
raises every such slope by at least 1, whereas \(P_{1,R}\) annihilates
inputs above its largest label. Hence the pair vanishes unless there
are labels \(\ell_L,\ell_R\) with \(\ell_R\ge\ell_L+1\). The same
support argument applies to all Jordan powers on these Riesz ranges.
For a surviving pair, the exponent from
\((\lambda_L\lambda_R)^m\) is \(N\) to the minus
\[
 3p^2(\ell_L^2+\ell_R^2)+6(e_L+e_R)-7/16 .
\]
Write the two slope lattices as
\[
 \ell_L=-\tfrac12+\tfrac89 i,\qquad
 \ell_R=\tfrac12+\tfrac89 j,\qquad i,j\in\mathbb Z.
\]
Admissibility is exactly \(j\ge i\), and the exponent becomes
\[
 \frac54+3(j-i)+\frac83(i^2+j^2)+6(e_L+e_R).
\]
Its unique minimum is therefore \(5/4\), at
\(\ell_L=-1/2,\ell_R=1/2,e_L=e_R=0\). Every other term has exponent
at least \(9/4\), since a different integer pair costs more than 1
and the least positive excitation is \(1/6\). Both individual
eigenvalues at the minimum are algebraically simple: the only
nearer-to-zero slope in each lattice has magnitude \(7/18\), whose
\(p^2\ell^2/2\) differs by \(1/18\), less than a positive excitation.
Other Jordan blocks add only polynomial powers in \(m\), uniformly
locally for the retained terms. Thus, with a sufficiently small radius
cut, the absolute normalized magnitude takes the form
\[
 N^{-5/4}|H(D)|+O(N^{-9/4+\eta'})\qquad(\eta'>0\text{ fixed}),
\]
where \(H(D)\) is the continuous leading coefficient (including the one-link offset), and errors are uniform locally.

It remains to prove that \(H\) is nonzero at some admissible width.
By \eqref{c4:Ang-4}, the ratio of the normalized marked and unmarked
magnitudes is
\[
 \frac{d_C(v)^2|w_-|}{d_C(h)}
                  \big|\tau^v_{N-1}/\tau_N\big|.
\]
The prefactor is a fixed positive constant. The determinant estimate
and \eqref{eq:one-mark-positive-input} therefore give a normalized
marked magnitude at least \(N^{-19/12+o(1)}\). If \(H\) vanished
throughout the admissible tail, the preceding error bound would instead
give \(O(N^{-9/4+\eta'})\). Choosing \(0<\eta'<2/3\) contradicts this
lower bound.

Fix an admissible \(D_*\) with \(H(D_*)\ne0\). For every sufficiently
large \(N\), round \(m\) so that the width
\(D_N=6\log N/m\) tends to \(D_*\). Continuity of \(H\) and the locally
uniform error then give normalized marked magnitude
\(N^{-5/4+o(1)}\) on the full sequence. Comparing with the unmarked
magnitude \(N^{-7/12+o(1)}\) through \eqref{c4:Ang-4} proves
\[
 \big|\tau^v_{N-1}/\tau_N\big|=N^{-2/3+o(1)} .
\]
